\documentclass{stp}

\newcommand{\en}[1]{\textit{#1}}
\newcommand{\todo}[1]{\textless{}TODO: #1\textgreater{}}
\newcommand{\pending}{[скриншот будет добавлен после настройки сервиса]}

\begin{document}

\stptitle{
  вид        = практическая,
  номер      = 2,
  тема       = {Деплой и мониторинг приложения},
  дисциплина = {ИТиВП},
  студент    = {А. С. Образцов},
  группа     = {320601},
  преподаватель = {Н. В. Хаджинова},
}

\section{Цель работы}

Целью работы является настройка автоматической сборки и развёртывания
серверной части сервиса планирования задач по методу \en{GTD}, её запуск на
бесплатной облачной платформе и подключение мониторинга: проверки
доступности, наблюдения за ресурсами, централизованных журналов и сбора
ошибок с уведомлениями. Основой служит \en{REST API} на \en{Node.js} и
\en{Express}, разработанный в предыдущих работах.

Приложение доступно по адресу \todo{продакшен-URL}, исходный код ---
в репозитории
\en{https://github.com/\allowbreak{}Alexandr-Obraztsov/\allowbreak{}gtd-task-service},
ветка \en{ipr22}. Ссылки на дашборды мониторинга и ошибок приведены в разделе
<<Мониторинг и деплой>> файла \en{README.md}; в черновике отчёта они остаются
заглушками, поскольку аккаунты внешних сервисов ещё не созданы.

\section{Ход работы}\label{sec:work}

\subsection{Выбор инструментов}

Инструменты выбирались по трём признакам: бесплатный тариф, доступность из
Беларуси и возможность настроить всё из репозитория, не прибегая к ручным
действиям при каждом выпуске. Выбранные средства сведены в таблице~\ref{tab:tools}.

\begin{stptable}
  \caption{Выбранные инструменты}\label{tab:tools}
  \begin{tabularx}{\linewidth}{|p{30mm}|p{36mm}|X|}\hline
    Задача & Инструмент & Обоснование\\\hline
    \en{CI/CD} & \en{GitHub Actions} & репозиторий уже на \en{GitHub},
    секреты и окружения встроены, достаточно бесплатных минут\\\hline
    Бэкенд & \en{Render} (\en{Docker}) & принимает \en{Dockerfile},
    запускает деплой по \en{Deploy Hook}, конфигурация хранится в
    \en{render.yaml}\\\hline
    Фронтенд & \en{Vercel} & деплой собранной статики одной командой
    \en{CLI}\\\hline
    База данных & \en{Neon} & управляемый \en{PostgreSQL} с \en{TLS} на
    бесплатном тарифе\\\hline
    Доступность & \en{UptimeRobot} & проверка раз в пять минут,
    публичная страница статуса, оповещения в \en{Telegram}\\\hline
    Журналы & \en{Better Stack} (\en{Logtail}) & подключается как ещё
    один транспорт \en{winston}, поиск по полям \en{JSON}\\\hline
    Ошибки & \en{Sentry} & стек вызова, окружение, релиз, группировка
    повторяющихся ошибок\\\hline
    Уведомления & \en{Telegram Bot API} & один запрос \en{curl} из
    \en{workflow}, без сторонних расширений\\\hline
  \end{tabularx}
\end{stptable}

\subsection{Проверка качества кода: \en{ESLint} и тесты}

Для статического анализа добавлен \en{ESLint} с конфигурацией
\en{eslint.config.js} (рекомендованные правила, запрет \en{console},
строгое сравнение, \en{prefer-const}) и скрипт \en{npm run lint}. Прогон
линтера по всему проекту завершился без замечаний, результат представлен на
рисунке~\ref{fig:lint}.

\stpimage[100mm]{img/lint.png}{Результат \en{npm run lint}}{fig:lint}

Автоматические тесты написаны на \en{Jest} и \en{supertest}. Они
запускают настоящее приложение (\en{createApp()}) с реальной базой
данных, а не заглушки. Проверяются: \en{/health} и \en{/health/ready},
ответ 401 без токена и с поддельным \en{JWT}, валидация (слабый пароль,
лишнее поле, битый \en{JSON}, неверный \en{Content-Type}), доступ к
\en{/metrics} для разных ролей, счётчик запросов, политика \en{CSP}
страницы \en{/dashboard} и маршрут тестовой ошибки. Всего 40 тестов в шести
файлах. Список тестов представлен на рисунках~\ref{fig:test1} и
\ref{fig:test2}.

\stpimage[120mm]{img/test1.png}{Выполнение тестов (начало)}{fig:test1}

\stpimage[120mm]{img/test2.png}{Выполнение тестов (окончание)}{fig:test2}

\subsection{Контейнеризация}

Образ описан в \en{Dockerfile} из двух стадий: в первой ставятся только
production"=зависимости (\en{npm ci --omit=dev} с кэшем \en{npm}), во
второй копируются исходники, и процесс запускается от пользователя
\en{node}, а не \en{root}. Файл \en{.dockerignore} исключает из контекста
\en{.env}, \en{node\_modules}, журналы, тесты и отчёты, поэтому секреты не
могут попасть в образ. Проверка живости встроена в образ инструкцией
\en{HEALTHCHECK}:

\begin{stplisting}
HEALTHCHECK --interval=30s --timeout=5s --start-period=20s \
    --retries=3 \
    CMD node -e "fetch('http://127.0.0.1:'+(process.env.PORT\
||10000)+'/health').then(r=>process.exit(r.ok?0:1))\
.catch(()=>process.exit(1))"
\end{stplisting}

Образ собран локально и запущен против контейнера с базой данных:
состояние контейнера стало \en{healthy}, а \en{/health/ready} вернул
\en{200}. Размер образа составил 417~МБ. Вывод команд представлен на
рисунке~\ref{fig:docker}.

\stpimage[140mm]{img/docker.png}{Сборка и запуск образа \en{Docker}}{fig:docker}

\subsection{\en{Workflow} \en{GitHub Actions}}

Файл \en{.github/workflows/deploy.yml} запускается при \en{push} в ветку
\en{main}, при \en{pull request} и вручную. Для \en{pull request}
выполняются только проверки. Начало файла с заданиями \en{test} и
\en{docker} представлено на рисунке~\ref{fig:wf1}. Задание \en{test}
поднимает \en{PostgreSQL~16} как \en{service container} с проверкой
готовности, ставит зависимости, запускает линтер и тесты. Задание
\en{docker} собирает образ через \en{buildx} с кэшем слоёв
\en{type=gha}, не публикуя его.

\stpimage[120mm]{img/workflow1.png}{Файл \en{workflow}: триггеры и задания \en{test}, \en{docker}}{fig:wf1}

Задания деплоя и уведомления представлены на рисунке~\ref{fig:wf2}.
Деплой бэкенда выполняется только после успеха проверок: запрос
\en{POST} на адрес \en{Deploy Hook} из секрета, затем опрос \en{/health}
до десяти минут. Деплой фронтенда на \en{Vercel} условный: он идёт, лишь
если в репозитории есть \en{client/package.json} и заданы секреты
\en{VERCEL\_*}. Оба деплоя и уведомление сами пропускаются, если нужных
секретов нет, поэтому \en{workflow} не падает в форке без настроенных
сервисов.

\stpimage[130mm]{img/workflow2.png}{Файл \en{workflow}: задания деплоя}{fig:wf2}

Страница запуска \en{workflow} в разделе \en{Actions} репозитория:
\pending.

\subsection{Секреты и переменные}

Значения секретов в репозитории не хранятся. Для \en{workflow} требуются
секреты \en{RENDER\_DEPLOY\_HOOK\_URL}, \en{VERCEL\_TOKEN},
\en{VERCEL\_ORG\_ID}, \en{VERCEL\_PROJECT\_ID}, \en{TELEGRAM\_BOT\_TOKEN},
\en{TELEGRAM\_CHAT\_ID} и переменная \en{PRODUCTION\_URL}; их назначение и
источник получения описаны в \en{README.md}. Секреты самого приложения
(\en{DATABASE\_URL}, \en{JWT\_SECRET}, \en{JWT\_REFRESH\_SECRET},
\en{METRICS\_TOKEN}, \en{SENTRY\_DSN}, \en{LOGTAIL\_SOURCE\_TOKEN})
вводятся только в панели \en{Render}: в \en{render.yaml} они помечены
\en{sync: false}, чтобы значения не попали в репозиторий.

\subsection{Деплой на \en{Render} и \en{Vercel}}

Конфигурация \en{Render} хранится в \en{render.yaml}: сервис типа \en{web}
с рантаймом \en{docker}, бесплатный план, регион \en{frankfurt},
\en{healthCheckPath: /health}. Автодеплой выключен
(\en{autoDeploy: false}), чтобы сборку запускал только \en{workflow} после
успешных проверок. База данных вынесена в \en{Neon}; приложение
подключается к ней по \en{TLS} (переменная \en{DATABASE\_SSL=true}, проверка
сертификата включена). Фрагмент конфигурации:

\begin{stplisting}
services:
  - type: web
    name: gtd-task-service
    runtime: docker
    plan: free
    region: frankfurt
    autoDeploy: false
    healthCheckPath: /health
    envVars:
      - key: NODE_ENV
        value: production
      - key: DATABASE_URL
        sync: false
\end{stplisting}

Для фронтенда добавлен \en{vercel.json}: команды установки и сборки
направлены в каталог \en{client}, выходной каталог --- \en{client/dist},
все пути, кроме \en{/assets/}, перенаправляются на \en{index.html}, чтобы
работала маршрутизация одностраничного приложения, добавлены заголовки
безопасности. Страница деплоя на платформе \en{Render}: \pending.
Страница проекта на \en{Vercel}: \pending.

\subsection{Мониторинг доступности}

Маршрут \en{GET /health} служит проверкой живости: он отвечает
\en{\{"status":"ok"\}} без обращения к базе данных, поэтому медленная
или недоступная база не приводит к перезапуску контейнера. Новый маршрут
\en{GET /health/ready} проверяет готовность: выполняет \en{SELECT 1} с
таймаутом 3~с и возвращает \en{200} либо \en{503}. Причина сбоя драйвера
наружу не отдаётся, поскольку текст ошибки может содержать адрес базы.
Ответы обоих маршрутов, а также отказ \en{/metrics} без токена представлены на
рисунке~\ref{fig:curl1}.

\stpimage[120mm]{img/curl1.png}{Ответы \en{/health}, \en{/health/ready} и \en{/metrics} без токена}{fig:curl1}

Для \en{UptimeRobot} в \en{README.md} приведена инструкция: монитор
\en{HTTP(s)} на \en{/health} с интервалом пять минут и ключевым словом
\en{"status":"ok"}, канал оповещений --- \en{Telegram}, публичная страница
статуса. Монитор на \en{/health/ready} предлагается ставить с большим
интервалом, потому что он обращается к базе и не даёт ей засыпать. Дашборд
\en{UptimeRobot}: \pending.

\subsection{Метрики и дашборд}

Промежуточный обработчик \en{metricsCollector} считает запросы: общее число,
находящиеся в обработке, разбивку по классам статуса и среднее время ответа.
Маршрут \en{GET /metrics} возвращает также время работы процесса,
\en{process.memoryUsage()} в мегабайтах, среднюю нагрузку и свободную
память системы. Ответ для запроса с токеном представлен на
рисунке~\ref{fig:curl2}.

\stpimage[85mm]{img/curl2.png}{Ответ \en{GET /metrics}}{fig:curl2}

Метрики раскрывают версию \en{Node.js}, объём памяти, нагрузку и профиль
трафика, то есть материал для разведки, поэтому доступ ограничен. Он открыт
по статическому токену \en{METRICS\_TOKEN} (его выдают сборщику метрик или
дашборду) либо по \en{JWT} пользователя с ролью \en{admin}. Токен
сравнивается через \en{crypto.timingSafeEqual} по хешам \en{SHA-256}, чтобы
время ответа не выдавало совпавшие символы; неверные попытки попадают в
журнал безопасности и ограничены лимитером. Маршрут \en{/health} оставлен
открытым намеренно: он нужен внешнему монитору и не содержит данных.

Страница \en{/dashboard} обновляется автоматически (5, 10 или 30~секунд).
Приложение отдаёт заголовок \en{CSP default-src 'none'}, а валидатор требует
пустые \en{params}, \en{query} и \en{body} у маршрутов без схемы.
Поэтому страница получила собственную политику (\en{script-src 'self'},
\en{connect-src 'self'}, без \en{unsafe-inline}): скрипт и стили вынесены в
отдельные файлы, данные выводятся через \en{textContent}, а не через
\en{innerHTML}. Токен вводится в форму и хранится в \en{sessionStorage}
вкладки, в адрес он не попадает. Обращения к самим метрикам в счётчик не
входят, иначе автообновление искажало бы статистику. Дашборд после ввода
токена, на фоне тестовой нагрузки, представлен на
рисунке~\ref{fig:dashboard}.

\stpimage[130mm]{img/dashboard.jpg}{Страница \en{/dashboard} в браузере}{fig:dashboard}

\subsection{Журналирование}

\en{Winston} уже писал журналы в консоль и в файлы с ротацией. Добавлен
транспорт \en{@logtail/winston}, который включается только при наличии
\en{LOGTAIL\_SOURCE\_TOKEN}; без токена поведение прежнее:

\begin{stplisting}
const logtailClient = config.logtail.sourceToken
  ? new Logtail(config.logtail.sourceToken,
      config.logtail.endpoint
        ? { endpoint: config.logtail.endpoint } : {})
  : null;

function remoteTransports() {
  return logtailClient
    ? [new LogtailTransport(logtailClient, { level: 'info' })]
    : [];
}
\end{stplisting}

К каждой записи добавлены поля \en{service} и \en{release}. Фиксируются
запуск сервера (порт, окружение, включённые \en{Sentry}, удалённые журналы и
метрики), каждый запрос (для кодов от 400 уровень \en{warn}), отклонённые
запросы с перечнем полей, не прошедших валидацию (значения не пишутся),
события безопасности и ошибки со стеком. Необработанные \en{uncaughtException}
и \en{unhandledRejection} записываются, отправляются в \en{Sentry}, после чего
процесс завершается, а платформа перезапускает контейнер. Локальные записи
после тестовой нагрузки представлены на рисунке~\ref{fig:logs}.

\stpimage[140mm]{img/logs.png}{Локальные журналы приложения}{fig:logs}

Интерфейс \en{Better Stack} с поступившими записями: \pending.

\subsection{Сбор ошибок: \en{Sentry}}

Модуль \en{src/utils/sentry.js} инициализирует \en{@sentry/node} только при
наличии \en{SENTRY\_DSN}. Перед отправкой событие очищается: удаляются
заголовки \en{Authorization} и \en{Cookie}, тело запроса и адрес пользователя,
остаётся только его \en{id}:

\begin{stplisting}
function scrubEvent(event) {
  if (event.request) {
    for (const name of Object.keys(event.request.headers || {})) {
      if (SENSITIVE_HEADERS.includes(name.toLowerCase())) {
        delete event.request.headers[name];
      }
    }
    delete event.request.cookies;
    delete event.request.data;
  }
  if (event.user) event.user = { id: event.user.id };
  return event;
}
\end{stplisting}

Обработчик ошибок приложения передаёт в \en{Sentry} все ошибки со статусом
от 500, а клиентские ошибки 4xx не отправляет, чтобы не засорять список.
Для проверки интеграции добавлен маршрут \en{POST /debug/sentry-test}. Он
безопасен: доступен только роли \en{admin}, только методом \en{POST}, не
принимает данных, а клиенту возвращает обычный ответ 500 без стека. Маршрут
оставлен и в \en{production}, иначе подключение нельзя проверить на живом
сервисе.

Внешнего аккаунта \en{Sentry} нет, поэтому интеграция проверена на локальном
приёмнике: значение \en{SENTRY\_DSN} указывало на заглушку
\en{127.0.0.1}, которая принимала запросы. Событие с тестовой ошибкой дошло до
приёмника с окружением, релизом и идентификатором пользователя, без
заголовка \en{Authorization}; журналы \en{@logtail/winston} дошли с
заголовком \en{Bearer}. Вывод проверки представлен на
рисунке~\ref{fig:sentry}. Это проверка кода, а не снимок настоящего сервиса.

\stpimage[140mm]{img/sentry_local.png}{Проверка отправки событий на локальный приёмник}{fig:sentry}

Дашборд \en{Sentry} с примером тестовой ошибки: \pending.

Для \en{React}"=клиента (ветка с каталогом \en{client}) в \en{README.md}
описана установка \en{@sentry/react}: вызов \en{Sentry.init} до рендера,
компонент \en{Sentry.ErrorBoundary} вокруг корня приложения, отдельный
\en{DSN} в переменной \en{VITE\_SENTRY\_DSN} и загрузка карт исходного кода
плагином \en{@sentry/vite-plugin}. В этой работе клиентский код не менялся.

\subsection{Уведомления в \en{Telegram}}

Последнее задание \en{workflow} (\en{notify}) выполняется всегда, даже если
предыдущие задания упали, определяет итог по результатам всех заданий и
отправляет сообщение запросом \en{curl} к \en{Bot API}. Данные, которые могут
содержать произвольные символы (имя ветки, автор), передаются через переменные
окружения, а не подставляются в текст команды, что исключает внедрение команд:

\begin{stplisting}
API=https://api.telegram.org/bot${TELEGRAM_BOT_TOKEN}
curl --fail --silent --show-error --max-time 30 \
  "${API}/sendMessage" \
  --data-urlencode "chat_id=${TELEGRAM_CHAT_ID}" \
  --data-urlencode "text=${TEXT}" > /dev/null
\end{stplisting}

Итог (успех, провал или отмена) определяется оператором \en{case} по
строке результатов всех заданий (\en{needs.*.result}).

Шаг пропускается, если секреты бота не заданы. Пример полученного
уведомления: \pending.

\section{Ответы на контрольные вопросы}

Список контрольных вопросов составлен автором работы самостоятельно по теме
задания.

\begin{stpnum}
  \item \textit{Что такое \en{CI/CD}?} \en{CI} (непрерывная интеграция) ---
    автоматическая сборка и проверка каждого изменения: линтер, тесты,
    сборка образа. \en{CD} (непрерывная доставка или развёртывание) ---
    автоматическая выкладка проверенной сборки в окружение. Вместе они
    сокращают время от коммита до работающей версии и убирают ручные,
    ошибкоопасные шаги.
  \item \textit{Как устроен \en{workflow} \en{GitHub Actions}?} Файл
    \en{YAML} в \en{.github/workflows} описывает события запуска
    (\en{push}, \en{pull\_request}, \en{workflow\_dispatch}) и задания
    (\en{jobs}). Каждое задание выполняется на отдельном исполнителе
    (\en{runner}) и состоит из шагов: готовых действий (\en{uses}) или
    команд (\en{run}). Порядок задаёт ключ \en{needs}, условия --- \en{if}.
  \item \textit{Как хранить секреты и почему не в репозитории?} В
    \en{GitHub Secrets} (шифруются, маскируются в журнале) и в панели
    платформы. Попавший в историю \en{git} секрет нельзя считать
    секретом: репозиторий публичный, история хранится вечно. Поэтому в
    репозитории только \en{.env.example}, а в \en{render.yaml}
    значения заменены на \en{sync: false}.
  \item \textit{Зачем в тестах \en{service container}?} Он поднимает
    настоящий \en{PostgreSQL} рядом с заданием, поэтому тесты проверяют
    реальные запросы и ограничения, а не заглушку. Контейнер создаётся
    чистым в каждом запуске, что исключает влияние данных прошлых прогонов.
  \item \textit{Зачем нужны \en{Docker}, многостадийная сборка и
    \en{.dockerignore}?} Образ даёт одинаковую среду на машине
    разработчика, в \en{CI} и на сервере. Многостадийная сборка оставляет
    в итоговом образе только нужное, запуск не от \en{root} снижает ущерб
    при взломе, а \en{.dockerignore} не пускает в образ секреты и мусор.
  \item \textit{Как кэш ускоряет сборку?} \en{actions/setup-node} с
    \en{cache: npm} сохраняет загруженные пакеты между запусками, а
    \en{buildx} с \en{type=gha} --- слои образа. Неизменившиеся слои
    (зависимости) берутся из кэша, пересобирается только код. Ключом кэша
    служат хеши \en{package-lock.json} и инструкций \en{Dockerfile}.
  \item \textit{Как работает деплой на \en{Render} и \en{Vercel}?}
    \en{Render} собирает образ из репозитория по \en{Dockerfile} и
    заменяет работающий экземпляр новым после успешной проверки
    \en{/health}; сборку запускает \en{Deploy Hook}, то есть секретный
    \en{URL}, на который \en{workflow} отправляет \en{POST}. \en{Vercel}
    принимает готовую сборку статики командой \en{vercel deploy --prebuilt}
    и раздаёт её через \en{CDN}.
  \item \textit{Чем проверка живости отличается от проверки готовности?}
    Живость (\en{/health}) отвечает, работает ли процесс; при провале
    платформа перезапускает контейнер. Готовность (\en{/health/ready})
    отвечает, может ли сервис обслуживать запросы, то есть доступна ли база;
    при провале трафик временно не направляют. Если включить в проверку
    живости базу, её сбой вызовет бессмысленные перезапуски приложения.
  \item \textit{Как работает внешний монитор \en{UptimeRobot}?} Раз в пять
    минут он запрашивает адрес из внешней сети, проверяет код ответа и
    ключевое слово, при сбое шлёт оповещение (в \en{Telegram}) и копит
    статистику доступности для публичной страницы. Ложные срабатывания
    уменьшает повторная проверка. Внешний монитор видит то, что видит
    пользователь, в отличие от проверки изнутри сети.
  \item \textit{Зачем нужны централизованные журналы?} Файловая система
    контейнера на \en{Render} временная, и локальные файлы пропадают при
    перезапуске. Внешний сервис хранит записи, позволяет искать по полям,
    фильтровать по уровню и строить оповещения. В журналы нельзя писать
    пароли, токены и значения полей запроса; поэтому при ошибках валидации
    пишутся только имена полей.
  \item \textit{Чем \en{Sentry} отличается от журналов?} Журнал --- поток
    всех событий, \en{Sentry} --- каталог проблем: он группирует одинаковые
    ошибки, показывает стек, частоту, первую и последнюю дату, окружение и
    релиз, позволяет назначить ответственного и пометить исправленной.
    Данные перед отправкой очищаются (\en{beforeSend}).
  \item \textit{Как безопасно сделать эндпоинт для тестовой ошибки?}
    Ограничить доступ ролью \en{admin}, разрешить только \en{POST}, не
    принимать входных данных, не раскрывать причину клиенту и записывать
    обращение в журнал. Альтернативы --- включать маршрут только вне
    \en{production} или по флагу окружения, но тогда проверить интеграцию
    на живом сервисе нельзя.
  \item \textit{Какие метрики собираются и как они защищены?} Время работы,
    память процесса, нагрузка и память системы, число запросов по классам
    статуса и среднее время ответа. Доступ --- по \en{METRICS\_TOKEN} или
    роли \en{admin}, токен сравнивается за постоянное время, неверные
    попытки журналируются и ограничиваются. Для нескольких экземпляров
    счётчики в памяти потребуют внешнего хранилища, например
    \en{Prometheus}.
  \item \textit{Какие существуют стратегии отката?} Откат к предыдущему
    образу (на \en{Render} --- \en{Rollback} в истории деплоев, на
    \en{Vercel} --- \en{Promote} прошлого деплоя), \en{blue-green}
    (два окружения и переключение трафика) и \en{canary} (новая версия
    получает долю трафика). Откат кода безопасен, пока изменения схемы базы
    только добавляют столбцы; в проекте миграций нет, схема создаётся
    \en{sync()}.
  \item \textit{Какие ограничения у бесплатных тарифов?} \en{Render} усыпляет
    сервис через 15 минут без трафика, холодный старт занимает до минуты, часы
    ограничены; \en{Neon} отключает вычисления при простое и ограничивает
    объём; \en{Sentry}, \en{Better Stack} и \en{UptimeRobot} ограничивают
    события, объём журналов и число мониторов. Смягчение: пинг монитором,
    допустимая задержка первого ответа, отбрасывание излишних данных.
  \item \textit{Как устроены уведомления в \en{Telegram} и какие у них
    риски?} Бот создаётся у \en{@BotFather}, \en{workflow} вызывает метод
    \en{sendMessage} с токеном бота и идентификатором чата. Токен лежит в
    секретах; он входит в адрес запроса, поэтому его нельзя печатать в
    журнал, а утёкший токен перевыпускается. Текст формируется из переменных
    окружения во избежание внедрения команд.
  \item \textit{Как не раскрыть секреты в \en{CI} и образах?} Не печатать
    значения, передавать их через \en{env}, а не аргументы команд, ограничить
    права токена (\en{permissions: contents: read}), не копировать \en{.env} в
    образ, пользоваться окружениями с защитой, регулярно перевыпускать
    ключи. Для запросов из форка секреты \en{GitHub} не передаёт, поэтому
    шаги деплоя в нём пропускаются.
\end{stpnum}

\stpunnumbered{Заключение}

В ходе работы серверная часть сервиса планирования задач подготовлена к
автоматическому развёртыванию и наблюдению. Добавлены линтер и 40 тестов,
\en{Dockerfile} (образ собран и проверен локально), \en{workflow}
\en{GitHub Actions} с проверками, сборкой образа, деплоем на \en{Render},
условным деплоем фронтенда на \en{Vercel} и уведомлениями в \en{Telegram},
конфигурации \en{render.yaml} и \en{vercel.json}.

Для наблюдения за работой реализованы проверки \en{/health} и
\en{/health/ready}, защищённые \en{/metrics} и \en{/dashboard}, журналы с
внешним транспортом \en{Better Stack}, отправка ошибок в \en{Sentry} с
очисткой данных и безопасный маршрут тестовой ошибки. Всё, что можно проверить
без внешних аккаунтов, проверено запуском; интеграции с \en{Sentry} и
\en{Better Stack} проверены на локальном приёмнике.

Не выполнено (подробности по каждому шагу даны в разделе~\ref{sec:work}): создание аккаунтов и сервисов (\en{Render}, \en{Neon},
\en{Vercel}, \en{UptimeRobot}, \en{Sentry}, \en{Better Stack}, бот
\en{Telegram}), добавление секретов в \en{GitHub}, первый запуск
\en{workflow} и реальный деплой. Поэтому ссылка на работающее приложение и
снимки внешних сервисов оставлены заглушками и будут добавлены после
настройки.

\end{document}
